\section{Varasemad tööd} \label{seotudtood}



\subsection{Kaugseire meetodite rakendamine surnud puidu hindamisel}



Surnud puidu, sealhulgas seisvate surnud puude (ingl \textit{snags}) ja maapinnal asuva lamapuidu (ingl \textit{downed deadwood}) seire on metsade süsinikuvaru ja elurikkuse hindamisel üks põhikomponent \parencite{marchi_airborne_2018, jarron_detection_2021}. 
Traditsioonilised välitöödel põhinevad inventeerimismeetodid, nagu joontransektid või proovitükid, võimaldavad küll kõrget mõõtetäpsust, kuid on suure töömahukuse tõttu raskesti skaleeritavad üleriigilistele maa-aladele \parencite{pesonen_airborne_2008, lopes_queiroz_estimating_2020}.

Sellest tulenev vajadus kaugseirelahenduste järele seab aga tingimuse eristada selgelt seisva ja lamapuidu tuvastamist, kuna neile kehtivad erinevad füüsikalised piirangud. Seisva surnud puidu puhul on puu võra õhust pealtvaates sageli nähtav, mis võimaldab edukalt kasutada passiivset optilist kaugseiret, näiteks spektraalanalüüsi aerofotodelt või satelliitpiltidelt \parencite{zielewska-buttner_detection_2020}. Seevastu lamapuit asub enamasti metsas, rinnete all, kus elusa puistu võrastik varjab tüvesid passiivsete sensorite eest, muutes vältimatuks aktiivse laserskaneerimise (LiDAR) kasutamise, mille laserkiired suudavad tungida läbi võra vahede maapinnani \parencite{marchi_airborne_2018}. Kuigi aerolaserskaneerimine (ingl \textit{Airborne Laser Scanning}, ALS) on muutunud operatiivses metsaseires standardseks tööriistaks, tugineb see alati kompromissile ruumilise ulatuse ja andmete detailsuse vahel \parencite{marchi_airborne_2018}. Kui suure punktitihedusega maapealne (TLS) ja mehitamata õhusõidukitel põhinev laserskaneerimine (ULS) võimaldavad detailset objektipõhist analüüsi, siis riiklike andmete puhul muutub kriitiliseks võime eristada tegelikku signaali mürast \parencite{jarron_detection_2021, vakimies_makaavan_2025}.

\subsection{Andmekvaliteedi ja taimkatte struktuuri mõju}

Lamapuidu tuvastamise edukus sõltub otseselt punktipilve tihedusest. Joyce jt (2019) näitasid eksperimentaalselt, et lamapuidu tuvastamise tõenäosus suureneb märkimisväärselt kuni umbes $16~p/m^2$ tiheduseni, mille järel täheldatakse täpsuse kasvus platooefekti \parencite{joyce_detection_2019}. Käesolevas töös kasutatavad andmed ($>18~p/m^2$) asuvad just selles optimaalses vahemikus, pakkudes piisavat detailsust üksikute lamapuidu objektide geomeetriliseks eristamiseks üleriigilises skaalas.

Tegelik maapinnani jõudvate laserimpulsside arv sõltub lisaks andmestiku üldisele tihedusele oluliselt ka metsatüübist ja fenoloogilisest faasist. Varasemad uuringud on näidanud, et lehtedeta periood on lamapuidu seireks optimaalseim, kuna raagus võrastik võimaldab laserkiirte parema läbitungimise maapinnani \parencite{marchi_airborne_2018}. Siiski varieerub penetratsioonivõime puistute lõikes – näiteks tihedates kuusikutes on võrastiku katvus suurem kui männikutes, mis piirab oluliselt maapinna lähedale jõudvate mõõtmiste arvu ja seeläbi ka tüvede tuvastatavust. Lisaks puistu ja alustaimestiku struktuurile mõjutavad detekteerimist objekti enda omadused: suurema diameetriga ja maapinnast kõrgemal asuvad tüved on punktipilves paremini eristatavad, samas kui peenem ja tugevalt lagunenud lamapuit võib jääda varjatuks alustaimestiku või maastiku mikroreljeefi tõttu \parencite{heinaro_airborne_2021, shin_morphological_2024}.

\subsection{Analüüsimeetodite areng: statistilistest mudelitest masinõppeni}

LiDAR-andmete analüüsimeetodid on arenenud koos andmete kvaliteedi ja kättesaadavusega, liikudes statistilistelt mudelitelt objektipõhiste ja masinõppel põhinevate lähenemisteni. Madala punktitihedusega andmete puhul ($4$--$10~p/m^2$) on laialdaselt kasutatud alal põhinevaid meetodeid (ABA), mis võimaldavad prognoosida puidu üldmahtu, kuid ei taga üksikobjektide tuvastamist \parencite{pesonen_airborne_2008}. Üksikute tüvede leidmiseks nendest andmetest on rakendatud näiteks automaatset sobitamist (\textit{template matching}) taimkatte kõrgusmudelil (CHM), mis on olnud edukas suurte tüvede tuvastamisel, kuid jätnud märkamata peenema lamapuidu \parencite{lindberg_detection_2013}.

Masinõppe ja objektipõhise analüüsi (OBIA) kombineerimine on võimaldanud täpsust tõsta ka keskmise tihedusega ($10$--$20~p/m^2$) andmetes. Blanchard jt (2011) saavutasid reeglipõhise masinõppe ja OBIA abil $73\%$-lise korrektsuse \parencite{blanchard_object-based_2011}, samas kui Mücke jt (2013) näitasid, et intensiivsusandmete ja OBIA integreerimine tagab avatud aladel kuni $75\%$-lise täpsuse \parencite{mucke_comparison_2013}. Nyström jt (2014) kasutasid joonsegmentide analüüsi ja Random Forest algoritmi, suutes tuvastada $64\%$ lamapuidu mahust isegi $10~p/m^2$ tiheduse juures \parencite{nystrom_detection_2014}. Need uuringud kinnitavad, et madala tiheduse korral tugineb edukas tuvastamine peamiselt kuju-indikaatoritele ja suurte objektide geomeetrilisele eristumisele. 

\subsection{Sügavõppe rakendamine ja skaleeritavuse väljakutsed}

Viimastel aastatel on lamapuidu tuvastamisel üha enam kasutatud sügavõppe meetodeid, eelkõige semantilist segmenteerimist ja üksikobjektide segmenteerimist (instance segmentation). Semantiline segmenteerimine määrab igale pikslile või punktile klassikuuluvuse, kuid ei erista sama klassi kuuluvaid objekte ükshaaval. Üksikobjektide segmenteerimine võimaldab seevastu eraldada individuaalseid lamavaid tüvesid, mis on oluline eriti juhul, kui tüved on osaliselt kattuvad või paiknevad lähestikku \parencite{polewski_instance_2021}. Polewski jt \parencite{polewski_instance_2021} kasutasid semantilise segmenteerimise tõenäosuskaarte individuaalsete lamavate tüvede eraldamiseks ning saavutasid kuni 0,93 täpsuse ja 0,82 saagise. Reder jt \parencite{reder_detection_2022} rakendasid U-Net arhitektuuri UAV-ortofotodelt tuulemurru tüvede tuvastamiseks, saavutades testandmestikul F1-skoorid vahemikus 0,726--0,756. Dietenberger jt \parencite{dietenberger_accurate_2025} kasutasid väga kõrge ruumilise lahutusega UAV-SfM andmetest tuletatud canopy-free ortomosaiiki ning U-Net mudelit ResNet-34 backbone’iga; nende F1-skoorid jäid sõltuvalt hindamistasemest vahemikku 0,73--0,96.

Uuemad uuringud on hakanud rasternäitajatelt ja võrastiku kõrgusmudelitel põhinevatelt lähenemistelt liikuma otse 3D-punktipilve analüüsi poole, sest CHM-põhised meetodid kirjeldavad eeskätt võrastiku ülemist pinda ning võivad seetõttu jätta alarinde objekte või keerukama vertikaalse struktuuri osaliselt arvestamata \parencite{xiang_automated_2024}. Näiteks Wielgosz jt \parencite{wielgosz_point2treep2tframework_2023} kasutasid Point2Tree raamistikus PointNet++-põhist semantilist segmenteerimist, et klassifitseerida punktipilves maapinda, taimestikku, tüvesid ja jämedat lamapuitu, saavutades semantilise segmenteerimise üldiseks F1-skooriks 0,92. Risti või ülekattuvalt paiknevate lamapuude eristamise puhul sobib viidata Polewski jt \parencite{polewski_instance_2021} tööle, kus kasutati U-Neti-põhiseid tõenäosuskaarte ja aktiivsete mitmikkontuuride meetodit; nende tulemused ulatusid kuni täpsuseni 0,93 ja saagiseni 0,82. Mask R-CNN-i ja mAP = 0,85 näitajat ei tohiks siduda risti paiknevate lamapuudega; see vastab paremini Sani-Mohammed jt \parencite{sani-mohammed_instance_2022} tööle seisvate surnud puude segmenteerimisel. Otse punktipilvest analüüsimine on lamapuidu ja madala taimestiku eristamiseks teoreetiliselt tugevam, kuna mudel saab kasutada objektide ruumilist ja vertikaalset struktuuri, mitte üksnes kõrgust maapinnast.

Üheks oluliseks takistuseks on endiselt tihe alustaimestik ja maapinna taimestik, mis võivad lamavaid tüvesid varjata ning põhjustada nii tuvastamata jäänud objekte kui ka valepositiivseid tuvastusi \parencite{heinaro_airborne_2021,heinaro_evaluating_2023}. Samuti näitavad varasemad uuringud, et automaatne tuvastamine on usaldusväärsem suuremate lamavate tüvede puhul; näiteks Heinaro jt \parencite{heinaro_airborne_2021} tuvastasid 78\% suurimatest lamapuudest (DBH $>300$ mm), samas kui kõigist üle 100 mm DBH-ga puudest tuvastati 30\%. Seetõttu tuleks automaatse lamapuidu tuvastuse tulemusi tõlgendada ettevaatlikult, eriti väiksema läbimõõduga tüvede, lagunenud puidu ja tiheda alustaimestikuga loodusmetsade puhul \parencite{marchi_airborne_2018,heinaro_airborne_2021,heinaro_evaluating_2023}.
\subsection{Uurimislünk ja töö eesmärk Eesti kontekstis}

Eelnevast teaduskirjandusest joonistub välja vastuolu kõrge täpsusega meetodite ja nende praktilise rakendatavuse vahel. Enamik edukaid lamapuidu uuringuid põhinevad lokaalsetel ja väga tihedatel andmetel (ULS, TLS), mis ei ole üleriigiliselt kättesaadavad \parencite{heinaro_airborne_2021, dietenberger_accurate_2025}. 

Eesti kontekstis on seni puudunud üleriigiliselt skaleeritav metoodika, mis rakendaks sügavõpet spetsiifiliselt lamapuidu automaatseks tuvastamiseks riiklikest aerolaserskanneerimise andmetest. Kuigi sarnaseid Eesti ALS-andmeid on edukalt rakendatud kuivenduskraavide tuvastamiseks \parencite{virro_detection_2025}, esitab lamapuit tuvastamismudelitel suuremaid väljakutseid objektide väikese vertikaalse ulatuse ja maapinnamüra tõttu. Käesoleva töö eesmärk on täita see lünk, kombineerides optimaalse tihedusega kevadised LiDAR-andmed ja sügavõppe, et luua esimene automatiseeritud lahendus lamapuidu üle-eestiliseks hindamiseks.
